\makeatletter
\let\pd@OrigRequirePackage\RequirePackage
\def\pd@dvipsname{dvips}
\def\RequirePackage{\@ifnextchar[\pd@RP@i{\pd@RP@i[\@empty]}}
\def\pd@RP@i[#1]#2{\@ifnextchar[{\pd@RP@ii{#1}{#2}}{\pd@RP@ii{#1}{#2}[\@empty]}}
\def\pd@RP@ii#1#2[#3]{%
  \def\pd@rpname{#2}%
  \ifx\pd@rpname\pd@hyperrefname
    \def\pd@newopts{}%
    \def\pd@first{1}%
    \ifx#1\@empty\else
      \@for\pd@tempopt:=#1\do{%
        \ifx\pd@tempopt\pd@dvipsname\else
          \ifx\pd@first\@empty
            \edef\pd@newopts{\pd@newopts,\pd@tempopt}%
          \else
            \edef\pd@newopts{\pd@tempopt}%
            \def\pd@first{}%
          \fi
        \fi
      }%
    \fi
    \ifx#3\@empty
      \pd@OrigRequirePackage[\pd@newopts]{#2}%
    \else
      \pd@OrigRequirePackage[\pd@newopts]{#2}[#3]%
    \fi
  \else
    \ifx#3\@empty
      \pd@OrigRequirePackage[#1]{#2}%
    \else
      \pd@OrigRequirePackage[#1]{#2}[#3]%
    \fi
  \fi
}
\def\pd@hyperrefname{hyperref}
\def\c@lor@to@ps#1 #2\@@{}%
\def\pdfmark#1{}%
\makeatother
\documentclass[
  style=simple,
  size=11pt,
  paper=screen,
  orient=landscape
]{powerdot}

\usepackage[utf8]{inputenc}
\usepackage{amsmath,amssymb}
\usepackage{booktabs}
\usepackage{colortbl}
\usepackage{tabularx}
\usepackage{fontawesome5}
\usepackage[most]{tcolorbox}
\usepackage{tikz}
\usetikzlibrary{positioning,calc,shapes.geometric,backgrounds}

% Color definitions (no trailing spaces)
\definecolor{monoBlack}{RGB}{18,18,18}
\definecolor{monoDark}{RGB}{35,39,47}
\definecolor{monoMuted}{RGB}{90,95,103}
\definecolor{monoLight}{RGB}{246,248,250}
\definecolor{monoBorder}{RGB}{210,215,222}

% TColorBox custom styles
\tcbset{
  monobox/.style={
    colback=monoLight,
    colframe=monoBlack,
    boxrule=0.8pt,
    arc=2pt,
    left=6pt,
    right=6pt,
    top=5pt,
    bottom=5pt,
    fonttitle=\bfseries\color{white},
    coltitle=white,
    colbacktitle=monoBlack
  },
  cardbox/.style={
    colback=white,
    colframe=monoBorder,
    boxrule=0.6pt,
    arc=2pt,
    left=6pt,
    right=6pt,
    top=4pt,
    bottom=4pt
  }
}

% Global Powerdot Setup
\pdsetup{
  lf={Meridian Systems Lab},
  cf={High-Scale Event-Driven Architecture},
  rf={\arabic{slide} / \pageref*{lastslide}}
}

\title{High-Scale Event-Driven Architecture\\[0.4em]\normalsize\mdseries Design Patterns for Zero-Downtime Financial Processing}
\author{Dr. Aris Thorne \and Elena Rostova\\[0.3em]\small Meridian Architecture \& Systems Lab}
\date{\today}

\begin{document}

% -----------------------------------------------------------------------------
% SLIDE 1: Title
% -----------------------------------------------------------------------------
\maketitle

% -----------------------------------------------------------------------------
% SLIDE 2: Agenda
% -----------------------------------------------------------------------------
\begin{slide}{Agenda \& Presentation Structure}
\begin{itemize}
  \item \textbf{System Foundations \& SLA Guarantees}: Core objectives, throughput targets, and sub-10ms p99 requirements.
  \item \textbf{Event Sourcing \& CQRS Pattern}: Decoupling state mutations from read-side query materialization.
  \item \textbf{Consensus \& Storage Architecture}: Raft-backed metadata consensus and partitioned event stream stores.
  \item \textbf{Topology \& Stream Engine}: End-to-end pipeline design from regional gateways to workers.
  \item \textbf{Resilience \& Operating Patterns}: Backpressure mechanisms, circuit breaking, and zero-downtime rollouts.
  \item \textbf{Architectural Takeaways}: Proven engineering guidelines for mission-critical financial systems.
\end{itemize}
\end{slide}

% -----------------------------------------------------------------------------
% SLIDE 3: Section Divider 1
% -----------------------------------------------------------------------------
\begin{wideslide}{Section 1: System Foundations \& SLA Guarantees}
\vspace{1.5em}
\begin{tcolorbox}[monobox,title=Architectural Core Focus]
Establishing the foundational pillars for high-throughput, low-latency, linearizable event processing across distributed financial nodes.
\end{tcolorbox}
\end{wideslide}

% -----------------------------------------------------------------------------
% SLIDE 4: Architectural Drivers
% -----------------------------------------------------------------------------
\begin{slide}{Architectural Drivers \& SLA Requirements}
\begin{minipage}{0.48\textwidth}
\begin{tcolorbox}[monobox,title=Throughput \& Latency SLAs]
\small
\begin{itemize}
  \item \textbf{Peak Volume}: 500,000 ops / sec continuous stream ingestion.
  \item \textbf{p99 Latency}: Sub-8ms end-to-end processing across regions.
  \item \textbf{Ingestion Gateway}: Distributed Nexa API gateway clusters with TLS termination.
\end{itemize}
\end{tcolorbox}
\end{minipage}%
\hfill
\begin{minipage}{0.48\textwidth}
\begin{tcolorbox}[monobox,title=Correctness \& Availability]
\small
\begin{itemize}
  \item \textbf{Consensus Model}: Strict linearizable state machine replication.
  \item \textbf{Uptime SLA}: 99.999\% availability with zero planned downtime.
  \item \textbf{Auditing}: Immutable cryptographic ledger append for all transactions.
\end{itemize}
\end{tcolorbox}
\end{minipage}
\end{slide}

% -----------------------------------------------------------------------------
% SLIDE 5: Event Sourcing & CQRS
% -----------------------------------------------------------------------------
\begin{slide}{Event Sourcing \& CQRS Pattern}
\begin{itemize}
  \item \textbf{Write-Path Isolation}: Commands mutate domain state strictly by appending immutable events to the stream store.
  \item \textbf{Optimistic Concurrency Control}: Sequence numbers validate state transitions without locking database tables.
  \item \textbf{Read-Path Projections}: Asynchronous projection workers consume events to populate read-optimized indexes.
  \item \textbf{Auditability \& Time Travel}: Complete historical sequence retention enables exact point-in-time state reconstruction.
\end{itemize}

\vspace{0.5em}
\begin{tcolorbox}[cardbox]
\centering
\small \textbf{Design Rule}: \textit{Never allow read-side queries to contend with write-side event log appends.}
\end{tcolorbox}
\end{slide}

% -----------------------------------------------------------------------------
% SLIDE 6: Section Divider 2
% -----------------------------------------------------------------------------
\begin{wideslide}{Section 2: Consensus \& Storage Architecture}
\vspace{1.5em}
\begin{tcolorbox}[monobox,title=Storage Engine Architecture]
Deep dive into Raft consensus, log partitioning, and comparative benchmark metrics across Meridian platform components.
\end{tcolorbox}
\end{wideslide}

% -----------------------------------------------------------------------------
% SLIDE 7: Raft Consensus State Machine
% -----------------------------------------------------------------------------
\begin{slide}{Raft-Backed Replicated State Machine}
\begin{minipage}{0.48\textwidth}
\begin{tcolorbox}[monobox,title=Consensus Mechanics]
\small
\begin{itemize}
  \item \textbf{Leader Election}: Heartbeat interval of 150ms with randomized election timeouts.
  \item \textbf{Log Replication}: Two-phase commit protocol guaranteeing quorum consensus before acknowledgment.
  \item \textbf{Membership Changes}: Single-server configuration changes to prevent split-brain states.
\end{itemize}
\end{tcolorbox}
\end{minipage}%
\hfill
\begin{minipage}{0.48\textwidth}
\begin{tcolorbox}[monobox,title=Log Compaction \& Snapshots]
\small
\begin{itemize}
  \item \textbf{Snapshot Trigger}: Automatic memory compaction every 100,000 committed entries.
  \item \textbf{Recovery Speed}: Fast startup initialization from state snapshot + recent log segment.
  \item \textbf{Storage Bounding}: Bounded disk memory consumption across node clusters.
\end{itemize}
\end{tcolorbox}
\end{minipage}
\end{slide}

% -----------------------------------------------------------------------------
% SLIDE 8: Storage Engine Benchmarks (Table Slide)
% -----------------------------------------------------------------------------
\begin{slide}{Storage Engine Benchmark Comparison}
\footnotesize
\centering
\begin{tabularx}{0.96\textwidth}{l X X c c}
\toprule
\textbf{Engine} & \textbf{Primary Use Case} & \textbf{Consensus Protocol} & \textbf{Read (p99)} & \textbf{Write (p99)} \\
\midrule
Nexa-Store & Immutable Event Ledger & Raft Quorum & 2.1 ms & 4.5 ms \\{}
Vertex-Log & Partitioned Stream Bus & Segment Lock & 1.4 ms & 3.1 ms \\{}
Nimbus-KV & Materialized Read Cache & Eventual Sync & 0.8 ms & 1.9 ms \\{}
Meridian-DB & Operational State Store & Multi-Paxos & 3.2 ms & 6.8 ms \\
\bottomrule
\end{tabularx}

\vspace{0.8em}
\begin{tcolorbox}[cardbox]
\small \textbf{Benchmark Summary}: Nexa-Store and Vertex-Log deliver consistent single-digit millisecond latency under 500k ops / sec load testing.
\end{tcolorbox}
\end{slide}

% -----------------------------------------------------------------------------
% SLIDE 9: System Topology Diagram (TikZ Slide)
% -----------------------------------------------------------------------------
\begin{slide}{End-to-End System Topology}
\centering
\vspace{-0.5em}
\begin{tikzpicture}[
  node distance=0.8cm and 0.9cm,
  box/.style={draw=monoBlack, thick, rectangle, rounded corners=3pt, fill=monoLight, inner sep=4pt, font=\scriptsize, align=center},
  arrow/.style={->, >=stealth, thick, draw=monoBlack}
]
  % Nodes
  \node[box] (client) {\faGlobe\\ \textbf{Client Apps}\\ (gRPC/HTTP)};
  \node[box, right=of client] (gw) {\faServer\\ \textbf{Nexa Gateway}\\ (TLS Layer)};
  \node[box, right=of gw] (bus) {\faProjectDiagram\\ \textbf{Vertex Bus}\\ (Event Stream)};
  \node[box, right=of bus] (worker) {\faCogs\\ \textbf{Stream Workers}\\ (Business Logic)};
  \node[box, below=0.8cm of worker] (store) {\faDatabase\\ \textbf{Nexa Store}\\ (Raft Ledger)};
  \node[box, left=0.9cm of store] (read) {\faLayerGroup\\ \textbf{Nimbus Cache}\\ (Read Models)};

  % Connections
  \draw[arrow] (client) -- (gw);
  \draw[arrow] (gw) -- (bus);
  \draw[arrow] (bus) -- (worker);
  \draw[arrow] (worker) -- (store);
  \draw[arrow] (worker) -- (read);
  \draw[arrow] (read) -- (gw);
\end{tikzpicture}

\vspace{0.2em}
\begin{tcolorbox}[cardbox]
\centering
\tiny \textbf{Data Flow}: Ingestion via Nexa Gateway $\rightarrow$ Vertex Stream Bus $\rightarrow$ Stream Workers $\rightarrow$ Nexa Store (Write) \& Nimbus Cache (Read).
\end{tcolorbox}
\end{slide}

% -----------------------------------------------------------------------------
% SLIDE 10: Section Divider 3
% -----------------------------------------------------------------------------
\begin{wideslide}{Section 3: Resilience \& Operating Patterns}
\vspace{1.5em}
\begin{tcolorbox}[monobox,title=Operational Reliability]
Implementation of reactive backpressure, automated circuit breaking, and canary deployment pipelines for zero-downtime operations.
\end{tcolorbox}
\end{wideslide}

% -----------------------------------------------------------------------------
% SLIDE 11: Fault Tolerance & Backpressure
% -----------------------------------------------------------------------------
\begin{slide}{Fault Tolerance \& Backpressure Strategies}
\begin{itemize}
  \item \textbf{Credit-Based Backpressure}: Downstream worker capacity dynamically signals ingestion gateways to adjust TCP window sizes.
  \item \textbf{Automated Circuit Breaking}: Triggers fallback routing if node error rate exceeds 0.05\% over a 10-second rolling window.
  \item \textbf{Dead Letter Queues (DLQ)}: Malformed poison-pill payloads are isolated to DLQs without stalling partition streams.
  \item \textbf{Self-Healing Partition Rebalance}: Failed cluster nodes trigger immediate partition reassignment within 3 seconds.
\end{itemize}
\end{slide}

% -----------------------------------------------------------------------------
% SLIDE 12: Zero-Downtime Deployment
% -----------------------------------------------------------------------------
\begin{slide}{Zero-Downtime Canary Deployment Pipeline}
\begin{minipage}{0.32\textwidth}
\begin{tcolorbox}[monobox,title=Phase 1: Shadow]
\small
\begin{itemize}
  \item 1\% traffic shadow.
  \item Compare response.
  \item Zero regression.
\end{itemize}
\end{tcolorbox}
\end{minipage}%
\hfill
\begin{minipage}{0.32\textwidth}
\begin{tcolorbox}[monobox,title=Phase 2: Canary]
\small
\begin{itemize}
  \item 25\% traffic shift.
  \item Monitor latency.
  \item Automated SLAs.
\end{itemize}
\end{tcolorbox}
\end{minipage}%
\hfill
\begin{minipage}{0.32\textwidth}
\begin{tcolorbox}[monobox,title=Phase 3: Rollout]
\small
\begin{itemize}
  \item 100\% promotion.
  \item DNS cutover.
  \item Auto-rollback.
\end{itemize}
\end{tcolorbox}
\end{minipage}
\end{slide}

% -----------------------------------------------------------------------------
% SLIDE 13: Scale Metrics
% -----------------------------------------------------------------------------
\begin{slide}{Production Scale \& Operational Metrics}
\begin{minipage}{0.48\textwidth}
\begin{tcolorbox}[cardbox]
\small
\textbf{Daily Ingestion Volume}\\[0.3em]
\Large \textbf{14.2 Billion Events / Day}\\[0.2em]
\tiny Continuous global processing across 8 regional data centers.
\end{tcolorbox}
\end{minipage}%
\hfill
\begin{minipage}{0.48\textwidth}
\begin{tcolorbox}[cardbox]
\small
\textbf{Event Replay Performance}\\[0.3em]
\Large \textbf{1,200,000 Events / sec}\\[0.2em]
\tiny Per-node recovery throughput during cold-start rebuilds.
\end{tcolorbox}
\end{minipage}

\vspace{0.8em}
\begin{minipage}{0.48\textwidth}
\begin{tcolorbox}[cardbox]
\small
\textbf{Recovery Time Objective (RTO)}\\[0.3em]
\Large \textbf{$<$ 5 Seconds}\\[0.2em]
\tiny Automatic failover and consensus recovery window.
\end{tcolorbox}
\end{minipage}%
\hfill
\begin{minipage}{0.48\textwidth}
\begin{tcolorbox}[cardbox]
\small
\textbf{Data Integrity SLA}\\[0.3em]
\Large \textbf{Zero Loss (RPO = 0)}\\[0.2em]
\tiny Guaranteed by Raft synchronous multi-region quorum write.
\end{tcolorbox}
\end{minipage}
\end{slide}

% -----------------------------------------------------------------------------
% SLIDE 14: Architectural Takeaways
% -----------------------------------------------------------------------------
\begin{slide}{Architectural Takeaways \& Engineering Guidelines}
\begin{itemize}
  \item \faCheck\ \textbf{Strict CQRS Separation}: Keep write-side state log append isolated from read-side projection serving.
  \item \faCheck\ \textbf{Bounded Snapshot Compaction}: Always bound event log recovery time with automated periodic snapshotting.
  \item \faCheck\ \textbf{Schema Evolution Contracts}: Enforce backward and forward schema compatibility across message schemas.
  \item \faCheck\ \textbf{End-to-End Idempotency}: Ensure consumers handle duplicate message delivery without state corruption.
\end{itemize}
\end{slide}

% -----------------------------------------------------------------------------
% SLIDE 15: Closing Slide & Q&A
% -----------------------------------------------------------------------------
\begin{wideslide}{Summary \& Engineering Q\&A}
\vspace{1.0em}
\begin{tcolorbox}[monobox,title=Meridian Architecture \& Engineering Lab]
\centering
\vspace{0.5em}
\Large \textbf{Thank You for Attending}\\[0.8em]
\small
\textbf{Dr. Aris Thorne} --- Principal Systems Architect \textbar\ \faEnvelope\ aris.thorne@meridian-systems.internal\\[0.3em]
\textbf{Elena Rostova} --- Lead Platform Engineer \textbar\ \faEnvelope\ elena.rostova@meridian-systems.internal\\[1.0em]
\tiny
Architecture Repository: \texttt{https://arch.meridian-systems.internal/docs/event-driven}
\vspace{0.5em}
\end{tcolorbox}
\end{wideslide}

\end{document}
